Er zijn twee soorten storage devices in een Unix-achtige omgeving zoals Linix. Er zijn de zogenaamde character en block devices. Character devices zijn over algemeen traag en verwerken data byte voor byte. Denk hierbij aan keyboards, muizen en geluidskaarten. Block devices verwerken vaak grote hoeveelheden data en de data wordt in grotere eenheden dan bytes verwerkt. Een harddisk kan bijvoorbeeld sectoren hebben van 512 bytes en dus wordt data per blok van 512 bytes verwerkt. Naast harddisks zijn ook, CD-ROMS, USB-sticks en Tape-drives block-devices.

In dit hoofdstuk behandelen we de block-devices.

